\documentclass[12pt,letterpaper,english]{amsart}
\usepackage[margin=1in]{geometry}
\usepackage{amsmath}   % utilities for mathematics
\usepackage{amsthm}    % utilities for theorem environment
\usepackage{amssymb}   % loads mathematical fonts
\usepackage[dvipsnames]{xcolor}
\usepackage{amscd}     % utilities for commutative diagrams

\usepackage{mathtools}

\usepackage{subcaption}
\usepackage{array}     % core package
\usepackage[sc]{mathpazo} % Palatino font with small caps
\usepackage[scr=boondoxo]{mathalfa}

\usepackage[unicode=true]{hyperref} % permits navigating on the pdf
\hypersetup{
    colorlinks=true,
    linkcolor=NavyBlue,
    citecolor=OliveGreen,
    urlcolor=blue,
    pdfencoding=auto
}
\usepackage{cleveref}  % smart referencing
\usepackage{caption}

\usepackage{graphics,graphicx} % permits putting images
\usepackage{tikz,tikz-cd}     % tool for diagrams
\usepackage{enumitem}

\usetikzlibrary{graphs,graphs.standard,calc}
\usetikzlibrary{decorations.pathreplacing}

\DeclareMathAlphabet{\mathsf}{OT1}{\sfdefault}{m}{n}

\newcommand{\nocontentsline}[3]{}
\newcommand{\tocless}[2]{\bgroup\let\addcontentsline=\nocontentsline#1{#2}\egroup}

\linespread{1.3}
\setlength\parindent{0.25in}
\usepackage{verbatim}
\usepackage{scalerel}
\usepackage{bbm}

\makeatletter
\@namedef{subjclassname@2020}{\textup{2020} Mathematics Subject Classification}
\makeatother


%%% -----------------------------------------------------------
%%% Theorem styles
%%% -----------------------------------------------------------
\newtheoremstyle{teoremas}%
{11pt}{11pt}{\itshape}{}{\bfseries}{}{.5em}{}

\theoremstyle{teoremas}
\newtheorem{theorem}{Theorem}[section]
\newtheorem{corollary}[theorem]{Corollary}
\newtheorem{lemma}[theorem]{Lemma}
\newtheorem{proposition}[theorem]{Proposition}

\newtheoremstyle{definition}%
{11pt}{11pt}{}{}{\bfseries}{}{.5em}{}

\theoremstyle{definition}
\newtheorem{definition}[theorem]{Definition}
\newtheorem{conjecture}[theorem]{Conjecture}
\newtheorem{example}[theorem]{Example}
\newtheorem{remark}[theorem]{Remark}

\crefname{theorem}{theorem}{theorems}
\Crefname{theorem}{Theorem}{Theorems}
\crefname{lemma}{lemma}{lemmas}
\Crefname{lemma}{Lemma}{Lemmas}
\crefname{proposition}{proposition}{propositions}
\Crefname{proposition}{Proposition}{Propositions}
\crefname{corollary}{corollary}{corollaries}
\Crefname{corollary}{Corollary}{Corollaries}


%%% -----------------------------------------------------------
%%% Operators and shortcuts
%%% -----------------------------------------------------------
\DeclareMathOperator{\rk}{rk}
\DeclareMathOperator{\crk}{crk}
\DeclareMathOperator{\cl}{cl}
\DeclareMathOperator{\PG}{PG}

\newcommand{\M}{\mathsf{M}}
\newcommand{\N}{\mathsf{N}}
\newcommand{\U}{\mathsf{U}}
\newcommand{\C}{\mathsf{C}}
\newcommand{\B}{\mathsf{B}}

\renewcommand{\H}{\mathrm{H}}
\newcommand{\uH}{\underline{\mathrm{H}}}
\newcommand{\Hstar}{\mathrm{H}^{*}}
\newcommand{\Gstar}{\mathrm{G}^{*}}
\newcommand{\PP}{\mathrm{P}}
\newcommand{\QQ}{\mathrm{Q}}
\newcommand{\ZZ}{\mathbb{Z}}
\newcommand{\FF}{\mathrm{F}}
\newcommand{\YY}{\mathrm{Y}}
\newcommand{\ZP}{\mathrm{Z}} % Z-polynomial

\newcommand{\rev}{\mathrm{rev}}
\newcommand{\hatQ}{\hat{\mathrm{Q}}}
\newcommand{\barchi}{\overline{\chi}}

\newcommand{\qbinom}[2]{\begin{bmatrix} #1 \\ #2 \end{bmatrix}_q}

\newcommand{\calL}{\mathcal{L}}
\newcommand{\calP}{\mathcal{P}}

\newcommand{\Se}{\mathcal{S}_e}
\newcommand{\Te}{\mathcal{T}_e}
\newcommand{\Ae}{\mathcal{A}_e}

\AtBeginDocument{%
   \def\MR#1{}
}

\allowdisplaybreaks


%%% -----------------------------------------------------------
%%% Bibliography
%%% -----------------------------------------------------------
%% Uncomment and point to your .bib file:
%\usepackage[backend=biber,style=alphabetic]{biblatex}
%\addbibresource{references.bib}


%%% -----------------------------------------------------------
%%% Title and metadata
%%% -----------------------------------------------------------
\title{Deletion Formulas}

\subjclass[2020]{Primary: 06A07, 05B35, 52B05, 05A20}


\begin{document}

\begin{abstract}
Aloha
\end{abstract}

\maketitle


%%% ============================================================
%%% SECTION 1 -- Introduction
%%% ============================================================
\section{Introduction}

%% TODO: Write the introduction.


%%% ============================================================
%%% SECTION 2 -- Preliminaries
%%% ============================================================
\section{Preliminaries}

\subsection{Conventions}
For a finite matroid~$\N$, write $E(\N)$ for its ground set,
$\mathcal{L}(\N)$ for its lattice of flats, and $\rk(\N)$ and
$\crk_{\N}(A) = \rk(\N) - \rk_{\N}(A)$ for rank and corank.
We use the standard minor notation~$\N|A$, $\N/A$, and~$\N \setminus e$.
If $a$ and~$b$ are functions on finite matroids, their \emph{convolution} is
\[
(a * b)_{\N}
= \sum_{A \in \mathcal{L}(\N)} a_{\N|A}\, b_{\N/A}.
\]
For a polynomial-valued invariant~$a$, its reversal is denoted by
\[
a^{\rev}_{\N}(x) = x^{\rk(\N)}\, a_{\N}(x^{-1}).
\]

The symbols $\PP$, $\H$, $\Hstar$, $\QQ$, $\YY$, $\ZP$, $\FF$,
and~$\Gstar$ denote the polynomial invariants introduced in the
applications below; their required convolution identities are recorded
there.  Whenever an invariant is extended by zero to matroids with loops,
that extension is stated explicitly at the point where it is used.

For a matroid~$\M$ and an element $e \in E(\M)$, define the
\emph{boundary sets}
\[
\Se(\M) = \{F \in \mathcal{L}(\M) \mid e \notin F
  \text{ and } F \cup \{e\} \in \mathcal{L}(\M)\},
\]
\[
\Te(\M) = \{B \in \mathcal{L}(\M) \mid e \in B,\;
  B \setminus \{e\} \notin \mathcal{L}(\M)\},
\]
and put $\overline{\Se}(\M) = \Se(\M) \setminus \{\emptyset\}$.
No non-coloop assumption is part of this notation; such an assumption
is made only when required by a deletion formula.

We define $\tau$ by
\[
\tau(\M) :=
\begin{cases}
[x^{(\rk(\M)-1)/2}]\, \PP_{\M}(x) & \text{if } \rk(\M) \text{ is odd},\\
0 & \text{if } \rk(\M) \text{ is even}.
\end{cases}
\]

\subsection{Flat calculus}
We recall the deletion formula for the characteristic polynomial:
\begin{lemma}\label{lem:char-poly-del}
Let $\N$ be a matroid and $e \in E(\N)$.
If $e$ is a coloop of~$\N$, then
\[
\chi_{\N}(x) = (x-1)\, \chi_{\N/e}(x).
\]
If $e$ is not a coloop of~$\N$, then
\[
\chi_{\N}(x) = \chi_{\N \setminus e}(x) - \chi_{\N/e}(x).
\]
\end{lemma}

The following two lemmas collect the flat decompositions and convolution
reindexings used throughout the deletion formulas.

\begin{lemma}[Flat calculus at a non-coloop]\label{lem:flat-calc}
Let $\M$ be loopless and let $e \in E(\M)$ be a non-coloop.  Put
\[
\Ae(\M) = \{A \in \mathcal{L}(\M) : e \notin A\},\quad
\mathcal{S} = \Se(\M),\quad
\mathcal{T} = \Te(\M).
\]
Then
\begin{align*}
\mathcal{L}(\M)
  &= \Ae(\M) \;\sqcup\; \{A \cup \{e\} : A \in \mathcal{S}\}
     \;\sqcup\; \mathcal{T}, \\
\mathcal{L}(\M \setminus e)
  &= \Ae(\M) \;\sqcup\; \{B \setminus \{e\} : B \in \mathcal{T}\}, \\
\mathcal{L}(\M/e)
  &= \mathcal{S} \;\sqcup\; \{B \setminus \{e\} : B \in \mathcal{T}\}.
\end{align*}
For every $A \in \Ae(\M)$, the matroid $\M/A$ is loopless,
$e$~is not a coloop of $\M/A$, and the maps $K \mapsto A \cup K$ identify
\[
\Se(\M/A) = \{F \setminus A : F \in \mathcal{S},\; A \subseteq F\},\quad
\Te(\M/A) = \{F \setminus A : F \in \mathcal{T},\; A \subseteq F\}.
\]
For every flat $B \in \mathcal{L}(\M)$ containing~$e$,
\[
\Se(\M|B) = \{F \in \mathcal{S} : F \subseteq B\},
\]
and, if $B \in \mathcal{T}$, then
\[
\Te(\M|B) = \{K \in \mathcal{T} : K \subseteq B\}.
\]
Finally, the map $(F,B) \mapsto (F, B \setminus F)$ is a bijection
\[
\{(F,B) : B \in \mathcal{T},\; F \in \Se(\M|B)\}
\;\longrightarrow\;
\{(F,K) : F \in \mathcal{S},\; K \in \Te(\M/F)\}.
\]
\end{lemma}

\begin{proof}
If a flat of~$\M$ contains~$e$, then either removing~$e$ leaves a flat,
which gives a unique member of~$\mathcal{S}$, or it does not, which gives
a member of~$\mathcal{T}$.  This proves the first decomposition.
If $C \in \mathcal{L}(\M \setminus e)$, then $\cl_{\M}(C)$ is either~$C$
or $C \cup \{e\}$.  In the former case $C \in \Ae(\M)$, while in the
latter $C = B \setminus \{e\}$ for the unique $B = \cl_{\M}(C) \in \mathcal{T}$;
the converse statements follow in the same way.  This proves the second
decomposition.  Finally,
\[
\cl_{\M/e}(C) = \cl_{\M}(C \cup \{e\}) \setminus \{e\}
\qquad (C \subseteq E(\M) \setminus \{e\}).
\]
Thus, if $C \in \mathcal{L}(\M/e)$ and $B = \cl_{\M}(C \cup \{e\})$,
then either $B \setminus \{e\} = C$ is a flat of~$\M$, so $C \in \mathcal{S}$,
or $B \in \mathcal{T}$ and $C = B \setminus \{e\}$.  Conversely, every flat
in either displayed class is a flat of~$\M/e$.
This proves the third decomposition.

If $A \in \Ae(\M)$, then $\M/A$ is loopless because $A$ is a flat, and
\[
\rk\bigl((\M/A) \setminus e\bigr) = \rk(\M \setminus e) - \rk(A)
= \rk(\M) - \rk(A) = \rk(\M/A),
\]
so $e$ is not a coloop of~$\M/A$.  The stated descriptions of its
boundary sets follow from the interval isomorphism between the flats
of~$\M/A$ and the flats of~$\M$ containing~$A$.  Likewise, the flats
of~$\M|B$ are the flats of~$\M$ contained in~$B$, which proves the two
restriction identities.  The last bijection has inverse
$(F,K) \mapsto (F, F \cup K)$; the flat descriptions just proved show
that these maps have the claimed domains.
\end{proof}

\begin{lemma}[Convolution reindexing]\label{lem:conv-reindex}
In the setting of \cref{lem:flat-calc}, let $u$ and~$v$ be functions on
finite matroids and put
\[
w_{\N} = \sum_{A \in \mathcal{L}(\N)} u_{\N|A}\, v_{\N/A}.
\]
Then
\begin{align}
w_{\M/e}
  &= \sum_{A \in \mathcal{S}} u_{\M|A}\, v_{\M/(A \cup \{e\})}
   + \sum_{B \in \mathcal{T}} u_{(\M|B)/e}\, v_{\M/B},
     \label{eq:reindex-1} \\
w_{\M|F}
  &= \sum_{\substack{A \in \mathcal{L}(\M) \\ A \subseteq F}}
     u_{\M|A}\, v_{(\M|F)/A}
     \qquad (F \in \mathcal{S}),
     \label{eq:reindex-2} \\
w_{\M/(L \cup \{e\})}
  &= \sum_{\substack{B \in \mathcal{L}(\M) \\ B \supseteq L \cup \{e\}}}
     u_{(\M|B)/(L \cup \{e\})}\, v_{\M/B}
     \qquad (L \in \mathcal{S}),
     \label{eq:reindex-3} \\
w_{\M/K}
  &= \sum_{\substack{B \in \mathcal{L}(\M) \\ B \supseteq K}}
     u_{(\M|B)/K}\, v_{\M/B}
     \qquad (K \in \mathcal{T}).
     \label{eq:reindex-4}
\end{align}
\end{lemma}

\begin{proof}
For the first identity, the third decomposition in \cref{lem:flat-calc}
applies, and the corresponding minors are
\[
(\M/e)|A = \M|A,\quad (\M/e)/A = \M/(A \cup \{e\})
\qquad (A \in \mathcal{S}),
\]
and
\[
(\M/e)|(B \setminus \{e\}) = (\M|B)/e,\quad
(\M/e)/(B \setminus \{e\}) = \M/B
\qquad (B \in \mathcal{T}).
\]
For the remaining identities, use the correspondence between flats
of~$\M/H$ and flats~$B$ of~$\M$ containing~$H$, together with
\[
(\M/H)|(B \setminus H) = (\M|B)/H,\quad
(\M/H)/(B \setminus H) = \M/B.
\]
Taking $H = F$, $L \cup \{e\}$, and~$K$, respectively, proves the claim.
\end{proof}


%%% ============================================================
%%% SECTION 3 -- Convolution Deletion Formula
%%% ============================================================
\section{Convolution Deletion Formula}

The next theorem records the bookkeeping needed to pass a pair of
deletion recurrences through a convolution.  It retains the terms
coming from contractions by nonflats, which vanish in the
loop-vanishing specializations used below.

\begin{theorem}[Convolution Deletion Formula]\label{thm:conv-del}
Let $R$ be a commutative ring, and let
$g$, $h$, $\Phi^h$, $\Psi^h$, $\Phi^g$, and~$\Psi^g$
be isomorphism-invariant functions on finite matroids with values in~$R$.
Define
\[
f_{\N} = \sum_{A \in \mathcal{L}(\N)} g_{\N|A}\, h_{\N/A}.
\]
Suppose that there are fixed scalars $\alpha_h, \alpha_g, \gamma_g \in R$
such that, for every pointed loopless matroid $(\N, p)$ with $p$~not a
coloop, one has
\begin{align*}
h_{\N} &= h_{\N \setminus p} + \alpha_h\, h_{\N/p}
  + \sum_{K_h \in \Se[p](\N)} \Phi^h\bigl(\N/(K_h \cup \{p\})\bigr)\, h_{\N|K_h}
  + \sum_{L_h \in \Te[p](\N)} \Psi^h(\N/L_h)\, h_{(\N|L_h)/p}, \\
g_{\N} &= g_{\N \setminus p} + \alpha_g\, g_{\N/p}
  + \sum_{K_g \in \Se[p](\N)} \Phi^g(\N|K_g)\, g_{\N/(K_g \cup \{p\})}
  + \sum_{L_g \in \Te[p](\N)} \Psi^g\bigl((\N|L_g)/p\bigr)\, g_{\N/L_g},
\end{align*}
and that, whenever $p$ is a coloop of a loopless matroid~$\N$,
\[
g_{\N} = \gamma_g\, g_{\N/p}
  + \sum_{K_g \in \Se[p](\N)} \Phi^g(\N|K_g)\, g_{\N/(K_g \cup \{p\})}.
\]
Let $\M$ be loopless and let $e \in E(\M)$ not be a coloop.
With
\[
\delta = \alpha_h + \gamma_g - \alpha_g,
\]
the following identity holds:
\begin{align*}
f_{\M} &= f_{\M \setminus e} + \alpha_g\, f_{\M/e}
  + \delta \sum_{A \in \Se(\M)} g_{\M|A}\, h_{\M/(A \cup \{e\})} \\
&\quad + \alpha_h \!\!\sum_{\substack{A \in \mathcal{L}(\M),\; e \notin A \\
    A \cup \{e\} \notin \mathcal{L}(\M)}}
  g_{\M|A}\, h_{\M/(A \cup \{e\})} \\
&\quad + \sum_{F \in \Se(\M)}
    \Phi^h\bigl(\M/(F \cup \{e\})\bigr)\, f_{\M|F} \\
&\quad + \sum_{F \in \Te(\M)} \Psi^h(\M/F)
    \sum_{\substack{A \in \mathcal{L}(\M|F) \\ e \notin A}}
    g_{(\M|F)|A}\, h_{(\M|F)/(A \cup \{e\})} \\
&\quad + \sum_{L \in \Se(\M)}
    \Phi^g(\M|L)\, f_{\M/(L \cup \{e\})} \\
&\quad + \sum_{K \in \Te(\M)}
    \Psi^g\bigl((\M|K)/e\bigr)\, f_{\M/K}.
\end{align*}
\end{theorem}

\begin{proof}
Put
\[
\mathcal{A} = \{A \in \mathcal{L}(\M) : e \notin A\},\quad
\mathcal{S} = \Se(\M),\quad
\mathcal{T} = \Te(\M).
\]
\cref{lem:flat-calc} gives all flat decompositions used below.
In particular, for every $A \in \mathcal{A}$, the matroid~$\M/A$ is
loopless and $e$~is not a coloop of~$\M/A$.  Set
\[
X = \sum_{A \in \mathcal{A}} g_{\M|A}\, h_{\M/A}.
\]
By the boundary-set description in \cref{lem:flat-calc}, the recurrence
for each~$h_{\M/A}$ is indexed by the flats $F \in \mathcal{S} \cup \mathcal{T}$
containing~$A$.  The relevant minor identities are
\[
(\M/A)/(K \cup \{e\}) = \M/(F \cup \{e\}),\quad
(\M/A)|K = (\M|F)/A,
\]
and
\[
(\M/A)/J = \M/F,\quad
((\M/A)|J)/e = ((\M|F)/A)/e,
\]
where $K = F \setminus A$ and $J = F \setminus A$.
Thus, after multiplying by~$g_{\M|A}$ and summing over $A \in \mathcal{A}$,
we group the deletion, contraction, $\Phi^h$-, and~$\Psi^h$-terms separately.
\cref{lem:flat-calc} and the restriction reindexing
identity~\eqref{eq:reindex-2} give
\begin{align}
X &= f_{\M \setminus e}
  - \sum_{B \in \mathcal{T}} g_{\M|(B \setminus \{e\})}\, h_{\M/B}
  \tag{deletion terms} \label{eq:del-terms} \\
&\quad + \alpha_h \sum_{A \in \mathcal{S}}
  g_{\M|A}\, h_{\M/(A \cup \{e\})}
  \tag{contraction terms with $A \cup \{e\}$ flat}
  \label{eq:contr-flat} \\
&\quad + \alpha_h \!\!\!\sum_{\substack{A \in \mathcal{L}(\M),\;
  e \notin A \\ A \cup \{e\} \notin \mathcal{L}(\M)}} \!\!\!
  g_{\M|A}\, h_{\M/(A \cup \{e\})}
  \label{eq:contr-nonflat} \\
&\quad + \sum_{F \in \mathcal{S}}
  \Phi^h\bigl(\M/(F \cup \{e\})\bigr)\, f_{\M|F}
  \tag{$\Phi^h$-boundary terms}
  \label{eq:Phih-terms} \\
&\quad + \sum_{F \in \mathcal{T}} \Psi^h(\M/F)
  \sum_{\substack{A \in \mathcal{L}(\M|F) \\ e \notin A}}
  g_{(\M|F)|A}\, h_{(\M|F)/(A \cup \{e\})}.
  \label{eq:Psih-terms}
\end{align}
The $\Phi^h$- and $\Psi^h$-boundary terms are obtained by exchanging the
sums over~$A$ and~$F$; for $F \in \mathcal{S}$, the resulting inner
convolution is~$f_{\M|F}$, by~\eqref{eq:reindex-2}.
Here the two contraction terms~\eqref{eq:contr-flat}
and~\eqref{eq:contr-nonflat} are separated according as $A \cup \{e\}$
is, or is not, a flat.  The deletion term follows from the deletion-flat
decomposition in \cref{lem:flat-calc}: if $B \in \mathcal{T}$, then $e$~is
a loop of~$\M/(B \setminus \{e\})$, and hence
$(\M \setminus e)/(B \setminus \{e\}) = \M/B$.

Put
\[
Y = \sum_{A \in \mathcal{S}} g_{\M|(A \cup \{e\})}\, h_{\M/(A \cup \{e\})}
  + \sum_{B \in \mathcal{T}} g_{\M|B}\, h_{\M/B}.
\]
Then $f_{\M} = X + Y$.
If $A \in \mathcal{S}$, then $\rk(A \cup \{e\}) = \rk(A) + 1$, so $e$ is
a coloop of $\M|(A \cup \{e\})$ and $(\M|(A \cup \{e\}))/e = \M|A$.
If $B \in \mathcal{T}$, then $e \in \cl_{\M}(B \setminus \{e\})$, so $e$
is not a coloop of~$\M|B$.  Moreover,
\[
\Se(\M|(A \cup \{e\})) = \{L \in \mathcal{S} : L \subseteq A\},
\]
and, for $B \in \mathcal{T}$,
\[
\Se(\M|B) = \{L \in \mathcal{S} : L \subseteq B\},\quad
\Te(\M|B) = \{K \in \mathcal{T} : K \subseteq B\}.
\]
Apply the same grouping argument to~$Y$, now using the recurrence for~$g$,
the preceding local boundary identities from \cref{lem:flat-calc},
and~\eqref{eq:reindex-3}--\eqref{eq:reindex-4}.  This yields
\begin{align}
Y &= \sum_{B \in \mathcal{T}} g_{\M|(B \setminus \{e\})}\, h_{\M/B}
  \notag \\
&\quad + \gamma_g \sum_{A \in \mathcal{S}} g_{\M|A}\, h_{\M/(A \cup \{e\})}
  + \alpha_g \sum_{B \in \mathcal{T}} g_{(\M|B)/e}\, h_{\M/B}
  \notag \\
&\quad + \sum_{L \in \mathcal{S}} \Phi^g(\M|L)\, f_{\M/(L \cup \{e\})}
  \notag \\
&\quad + \sum_{K \in \mathcal{T}}
  \Psi^g\bigl((\M|K)/e\bigr)\, f_{\M/K}.
  \label{eq:Y-expansion}
\end{align}

Finally, \eqref{eq:reindex-1}~gives
\[
f_{\M/e} = \sum_{A \in \mathcal{S}} g_{\M|A}\, h_{\M/(A \cup \{e\})}
  + \sum_{B \in \mathcal{T}} g_{(\M|B)/e}\, h_{\M/B}.
\]
Adding~\eqref{eq:del-terms} and~\eqref{eq:Y-expansion}, the two sums
indexed by~$\mathcal{T}$ with factor~$g_{\M|(B \setminus \{e\})}\, h_{\M/B}$
cancel.  The displayed identity for~$f_{\M/e}$ and the equality
$\delta = \alpha_h + \gamma_g - \alpha_g$ then transform the remaining
scalar terms into the first two lines of the asserted formula.
The four boundary sums are already the remaining four lines of that formula.
\end{proof}

\begin{corollary}[Loop-vanishing convolution deletion formula]
\label{cor:loop-vanish}
In the setting of \cref{thm:conv-del}, additionally assume that
$h_{\N} = 0$ whenever $\N$ has a loop, that $\Psi^h = 0$,
and that $\Phi^g(\emptyset) = 0$.  Then
\begin{align*}
f_{\M} &= f_{\M \setminus e} + \alpha_g\, f_{\M/e}
  + \delta \sum_{A \in \Se(\M)} g_{\M|A}\, h_{\M/(A \cup \{e\})} \\
&\quad + \sum_{F \in \Se(\M)}
    \Phi^h\bigl(\M/(F \cup \{e\})\bigr)\, f_{\M|F} \\
&\quad + \sum_{L \in \Se(\M)}
    \Phi^g(\M|L)\, f_{\M/(L \cup \{e\})} \\
&\quad + \sum_{K \in \Te(\M)}
    \Psi^g\bigl((\M|K)/e\bigr)\, f_{\M/K}.
\end{align*}
\end{corollary}

\begin{proof}
Let $A \in \mathcal{L}(\M)$ avoid~$e$ and suppose that $A \cup \{e\}$
is not a flat.  Choose
$z \in \cl_{\M}(A \cup \{e\}) \setminus (A \cup \{e\})$.
Then
\[
\rk_{\M/(A \cup \{e\})}(\{z\})
= \rk_{\M}(A \cup \{e,z\}) - \rk_{\M}(A \cup \{e\}) = 0,
\]
so $z$ is a loop of~$\M/(A \cup \{e\})$.
The loop-vanishing hypothesis therefore makes every summand in the
nonflat-boundary sum zero.  The hypothesis $\Psi^h = 0$ eliminates the
$\Psi^h$-sum, and $\Phi^g(\emptyset) = 0$ makes the $L = \emptyset$
summand vanish.  The theorem now reduces exactly to the displayed formula.
\end{proof}

\begin{remark}\label{rem:loop-vanish-needed}
The loop-vanishing hypothesis in the corollary cannot be omitted:
for example, $h_{\N} = 2^{|E(\N)|}$ and $g_{\N} = 1$, with
$\alpha_h = \gamma_g = 1$, $\alpha_g = 0$, and all boundary weights zero,
satisfy the recurrences but the nonflat-boundary sum for~$\U_{2,3}$ is nonzero.
\end{remark}


%%% ============================================================
%%% SECTION 4 -- Inverse Deletion Formula
%%% ============================================================
\section{Inverse Deletion Formula}

The next result is a companion to the Convolution Deletion Formula.
It makes precise how a deletion recurrence for a convolution kernel
transfers to its inverse.  The residual term records exactly what can
remain when contractions by nonflats create loops.

\begin{theorem}[Inverse Deletion Formula]\label{thm:inv-del}
Let $R$ be a commutative ring, and let $\kappa$, $\lambda$, and~$\phi$
be isomorphism-invariant functions on finite matroids with values in~$R$.
Suppose that
\[
\kappa_{\emptyset} = \sigma \in \{1, -1\}
\]
and that $\lambda$ is a right convolution inverse of~$\kappa$, namely
\begin{equation}\label{eq:conv-inv}
\sum_{F \in \mathcal{L}(\N)} \kappa_{\N|F}\, \lambda_{\N/F}
= \delta_{\emptyset, \N}
\quad \text{for every finite matroid } \N.
\end{equation}
Assume that there are fixed scalars $\alpha, \gamma, \varepsilon \in R$
such that, for every pointed loopless matroid $(\N, p)$, one has
\begin{alignat}{2}
\kappa_{\N} &= \kappa_{\N \setminus p} + \alpha\, \kappa_{\N/p}
  + \sum_{F \in \Se[p](\N)}
    \phi\bigl(\N/(F \cup \{p\})\bigr)\, \kappa_{\N|F}
  &\quad& \text{if $p$ is not a coloop},
  \label{eq:kappa-noncoloop} \\
\kappa_{\N} &= \gamma\, \kappa_{\N/p}
  + \varepsilon\, \mathbf{1}_{\{\N \cong \U_{1,1}\}}
  &\quad& \text{if $p$ is a coloop}.
  \label{eq:kappa-coloop}
\end{alignat}
Then, for every loopless matroid~$\M$ and every non-coloop
$e \in E(\M)$,
\begin{align}
\lambda_{\M} &= \lambda_{\M \setminus e}
  + \sigma(\alpha - \gamma)
    \sum_{A \in \Se(\M)} \kappa_{\M|A}\, \lambda_{\M/(A \cup \{e\})}
  - \sigma\varepsilon\, \mathbf{1}_{\{\emptyset \in \Se(\M)\}}\,
    \lambda_{\M/e}
    \notag \\
&\quad - \sigma \!\!\sum_{\substack{A \in \mathcal{L}(\M) \\
    e \notin A,\; A \neq \emptyset}} \!\!
    \kappa_{\M|A}\bigl(\lambda_{\M/A} - \lambda_{(\M/A) \setminus e}\bigr)
    \notag \\
&\quad - \sigma \sum_{A \in \Se(\M)} \kappa_{\M|A}
    \sum_{B \in \Te(\M/A)}
    \phi\bigl(((\M/A)|B)/e\bigr)\, \lambda_{(\M/A)/B}.
    \label{eq:inv-del}
\end{align}
\end{theorem}

\begin{proof}
Put
\[
\mathcal{A} = \{A \in \mathcal{L}(\M) : e \notin A\},\quad
\mathcal{S} = \Se(\M),\quad
\mathcal{T} = \Te(\M),\quad
\eta = \mathbf{1}_{\{\emptyset \in \mathcal{S}\}}.
\]
First suppose that $\sigma = 1$.
We use \cref{lem:flat-calc}.  In particular, a flat of~$\M$ containing~$e$
is either $A \cup \{e\}$ for a unique $A \in \mathcal{S}$, or belongs
to~$\mathcal{T}$.

Let
\[
X = \sum_{\substack{B \in \mathcal{L}(\M) \\ e \in B}}
  \kappa_{\M|B}\, \lambda_{\M/B}.
\]
For $A \in \mathcal{S}$, the element~$e$ is a coloop of
$\M|(A \cup \{e\})$, whereas for $B \in \mathcal{T}$ it is not a coloop
of~$\M|B$.
Applying~\eqref{eq:kappa-coloop} in the first case
and~\eqref{eq:kappa-noncoloop} in the second gives
\begin{align}
X &= \gamma \sum_{A \in \mathcal{S}}
    \kappa_{\M|A}\, \lambda_{\M/(A \cup \{e\})}
  + \sum_{B \in \mathcal{T}}
    \kappa_{\M|(B \setminus \{e\})}\, \lambda_{\M/B}
    \notag \\
&\quad + \alpha \sum_{B \in \mathcal{T}}
    \kappa_{(\M|B)/e}\, \lambda_{\M/B}
  + \sum_{A \in \mathcal{S}} \kappa_{\M|A}
    \sum_{C \in \Te(\M/A)}
    \phi\bigl(((\M/A)|C)/e\bigr)\, \lambda_{(\M/A)/C}
    \notag \\
&\quad + \eta \varepsilon\, \lambda_{\M/e}.
    \label{eq:X-inv}
\end{align}
Indeed, the exceptional summand in the coloop recurrence occurs only
when $A = \emptyset$, and this index is present precisely when
$\emptyset \in \mathcal{S}$.  For the last term, we have interchanged
the two boundary sums using the final bijection of \cref{lem:flat-calc}.

The proof then proceeds by induction on $|E(\M)|$.
Subtracting $X$ from the convolution identity~\eqref{eq:conv-inv}
and using~\eqref{eq:reindex-1} to identify~$f_{\M/e}$ yields the
result after collecting terms.
The case $\sigma = -1$ follows by replacing~$\kappa$ with~$-\kappa$.
\end{proof}

\begin{corollary}[Loop-vanishing inverse deletion formula]
\label{cor:inv-loop-vanish}
In the setting of \cref{thm:inv-del}, additionally assume that
$\kappa_{\N} = 0$ whenever $\N$ has a loop, and that $\phi(\N) = 0$
whenever $\N$ has a loop.
If $\varepsilon = 0$, or more generally if
$\varepsilon \neq 0$ only when $\N \cong \U_{1,1}$, then the last two lines
of~\eqref{eq:inv-del} vanish, and
\begin{align*}
\lambda_{\M} &= \lambda_{\M \setminus e}
  + \sigma(\alpha - \gamma) \sum_{A \in \Se(\M)}
    \kappa_{\M|A}\, \lambda_{\M/(A \cup \{e\})}
  - \sigma\varepsilon\, \mathbf{1}_{\{\emptyset \in \Se(\M)\}}\,
    \lambda_{\M/e} \\
&\quad + \sum_{B \in \Te(\M)}
    \phi\bigl((\M|B)/e\bigr)\, \lambda_{\M/B}.
\end{align*}
\end{corollary}


%%% ============================================================
%%% SECTION 5 -- Deletion Formulas
%%% ============================================================
\section{Deletion Formulas}

\subsection{The Kazhdan--Lusztig Polynomial}

\begin{theorem}[Kazhdan--Lusztig deletion formula]\label{thm:KL-del}
Let $\M$ be loopless and let $e \in E(\M)$ be a non-coloop.  Then
\[
\PP_{\M}(x) = \PP_{\M \setminus e}(x) - x\, \PP_{\M/e}(x)
  + \sum_{F \in \Se(\M)}
    \tau\bigl(\M/(F \cup \{e\})\bigr)\, x^{\crk_{\M}(F)/2}\,
    \PP_{\M|F}(x).
\]
\end{theorem}

\begin{proof}
Apply the three flat decompositions of \cref{lem:flat-calc} to the
convolution identity
$\PP^{\rev}_{\M}(x) = \sum_{B \in \mathcal{L}(\M)}
\chi_{\M|B}(x)\, \PP_{\M/B}(x)$.
Using \cref{lem:char-poly-del} on the factors indexed by flats containing~$e$ gives
\begin{align}
\PP^{\rev}_{\M}(x) &= \sum_{A \not\ni e} \chi_{\M|A}\, \PP_{\M/A}
  + \sum_{A \in \Se(\M)} (x-1)\, \chi_{\M|A}\, \PP_{\M/(A \cup \{e\})}
    \notag \\
&\quad + \sum_{B \in \Te(\M)}
    (\chi_{\M|B \setminus \{e\}} - \chi_{(\M|B)/e})\, \PP_{\M/B}.
    \label{eq:KL-expansion}
\end{align}
The flats of $\M \setminus e$ and $\M/e$ partition analogously, giving
\begin{align*}
\PP^{\rev}_{\M \setminus e}(x)
  &= \sum_{A \not\ni e} \chi_{\M|A}\, \PP_{(\M/A) \setminus e}
  + \sum_{B \in \Te(\M)}
    \chi_{\M|B \setminus \{e\}}\, \PP_{\M/B}, \\
\PP^{\rev}_{\M/e}(x)
  &= \sum_{A \in \Se(\M)} \chi_{\M|A}\, \PP_{\M/(A \cup \{e\})}
  + \sum_{B \in \Te(\M)}
    \chi_{(\M|B)/e}\, \PP_{\M/B}.
\end{align*}
The alternating sum
$\PP^{\rev}_{\M} - \PP^{\rev}_{\M \setminus e} + \PP^{\rev}_{\M/e}$
cancels terms corresponding to~$\Te(\M)$, yielding
\begin{equation}\label{eq:KL-alt-sum}
\PP^{\rev}_{\M} - \PP^{\rev}_{\M \setminus e} + \PP^{\rev}_{\M/e}
= \sum_{A \in \Se(\M) \cup \{\emptyset\}}
  \chi_{\M|A}\bigl(\PP_{\M/A} - \PP_{(\M/A) \setminus e}
  + x\, \PP_{(\M/A)/e}\bigr),
\end{equation}
where terms with $A \notin \Se(\M)$ vanish because $e$ becomes parallel
to another element in~$\M/A$, forcing
$\PP_{\M/A} = \PP_{(\M/A) \setminus e}$ and $\PP_{(\M/A)/e} = 0$.

For every loopless minor~$\N$ in which $e$ is a non-coloop, put
\[
C_e(\N) = \PP_{\N} - \PP_{\N \setminus e} + x\, \PP_{\N/e}.
\]
Since $C_e(\M)^{\rev} = x^{\rk(\M)}\, C_e(\M)(x^{-1})$,
equation~\eqref{eq:KL-alt-sum} simplifies to
\begin{equation}\label{eq:Ce-recursion}
C_e(\M)^{\rev} - C_e(\M)
= \sum_{\emptyset \neq A \in \Se(\M)}
  \chi_{\M|A}(x)\, C_e(\M/A).
\end{equation}

Define
\[
R_e(\N) = \sum_{F \in \Se(\N)}
  \tau\bigl(\N/(F \cup \{e\})\bigr)\, x^{\crk_{\N}(F)/2}\,
  \PP_{\N|F}(x).
\]
We claim, by induction on $|E(\M)|$, that $C_e(\M) = R_e(\M)$.
The defining convolution identity for~$\PP$ gives
\begin{align*}
R_e(\M)^{\rev}
&= \sum_{F \in \Se(\M)}
  \tau\bigl(\M/(F \cup \{e\})\bigr)\, x^{\crk_{\M}(F)/2}\,
  \PP^{\rev}_{\M|F}(x) \\
&= \sum_{F \in \Se(\M)}
  \tau\bigl(\M/(F \cup \{e\})\bigr)\, x^{\crk_{\M}(F)/2}
  \sum_{A \in \mathcal{L}(\M|F)} \chi_{\M|A}(x)\, \PP_{(\M|F)/A}(x) \\
&= R_e(\M) + \sum_{\emptyset \neq A \in \Se(\M)}
  \chi_{\M|A}(x)\, R_e(\M/A).
\end{align*}
Indeed, in the last equality we exchange the sums.
If $A \subseteq F$ with $F \in \Se(\M)$, then $A \in \Se(\M)$;
and the flats $F/A$ that arise are precisely the members
of~$\Se(\M/A)$.  The corresponding restrictions, contractions,
coranks, and $\tau$-coefficients agree.
This is also the point at which the formula is automatically zero
when $e$ has a parallel mate: then $\Se(\M) = \emptyset$.

By the induction hypothesis and~\eqref{eq:Ce-recursion}, the difference
$C_e(\M) - R_e(\M)$ is invariant under reversal.  It remains to use the
usual degree bound.  Every nonempty-flat summand of~$R_e(\M)$ has degree
strictly smaller than $\rk(\M)/2$.  The same holds for~$C_e(\M)$ unless
$\rk(\M)$ is even; in that case its coefficient in degree $\rk(\M)/2$
comes only from $x\, \PP_{\M/e}$ and equals
\[
[x^{(\rk(\M)-2)/2}]\, \PP_{\M/e}(x) = \tau(\M/e).
\]
This is exactly the degree-$\rk(\M)/2$ coefficient of the $F = \emptyset$
summand of~$R_e(\M)$ (and both coefficients are zero when
$\emptyset \notin \Se(\M)$).  Hence
$\deg(C_e(\M) - R_e(\M)) < \rk(\M)/2$.  A polynomial of degree strictly
less than half its reversal degree which is invariant under reversal
is zero.  Thus $C_e(\M) = R_e(\M)$, which is the asserted formula.
\end{proof}


\subsection{The Chow Polynomial}

\begin{corollary}[Chow polynomial deletion formula]\label{cor:chow-del}
Let $\M$ be a loopless matroid and let $e \in E(\M)$.  Then
\begin{equation}\label{eq:chow-del}
\H_{\M}(x) = \H_{\M \setminus e}(x)
  + x \sum_{G \in \overline{\Se}(\M)}
    \H_{\M|G}(x)\, \H_{\M/(G \cup \{e\})}(x).
\end{equation}
\end{corollary}

\begin{proof}
Use the defining expression for $\Se(\M)$ also when $e$ is a coloop.
Extend $\H$ by~$0$ to matroids with loops, with $\H_{\emptyset} = 1$,
and define
\[
\kappa_{\N} =
\begin{cases}
-1, & \N = \emptyset, \\
\dfrac{\chi_{\N}(x)}{x - 1}, & \N \neq \emptyset,
\end{cases}
\qquad \lambda_{\N} = -\H_{\N}(x).
\]
The defining convolution identity for~$\H$ says exactly that $\lambda$ is
the right convolution inverse of~$\kappa$.

By the characteristic-polynomial deletion identities, the hypotheses
of \cref{thm:inv-del} hold with
\[
\alpha = -1,\quad \gamma = x-1,\quad \varepsilon = x,\quad \phi = 0.
\]
We prove the result by induction on $|E(\M)|$; the one-element case is
immediate.  First suppose that $e$ is not a coloop and put
$D_e(\N) = \H_{\N} - \H_{\N \setminus e}$.
If $\emptyset \in \Se(\M)$, the $A = \emptyset$ summand in the inverse
deletion formula cancels the exceptional term~$x\, \H_{\M/e}$;
if $\emptyset \notin \Se(\M)$, both terms are absent.
Thus \cref{thm:inv-del} gives
\begin{equation}\label{eq:chow-induction}
(x-1)\, D_e(\M)
- \sum_{\substack{A \in \mathcal{L}(\M) \\
    e \notin A,\; A \neq \emptyset}}
  \chi_{\M|A}(x)\, D_e(\M/A)
= x \sum_{G \in \overline{\Se}(\M)}
  \chi_{\M|G}(x)\, \H_{\M/(G \cup \{e\})}(x).
\end{equation}
For every nonempty flat~$A$ avoiding~$e$, the induction hypothesis
applies to~$\M/A$.  Reindexing the resulting sums by $G = A \cup C$,
as in \cref{lem:flat-calc}, gives
\begin{align*}
\sum_{\substack{A \in \mathcal{L}(\M) \\
    e \notin A,\; A \neq \emptyset}}
  \chi_{\M|A}\, D_e(\M/A)
&= x \sum_{G \in \overline{\Se}(\M)}
  \Biggl(\sum_{\substack{A \in \mathcal{L}(\M|G) \\
    \emptyset \neq A \subsetneq G}}
  \chi_{\M|A}\, \H_{(\M|G)/A}\Biggr)
  \H_{\M/(G \cup \{e\})} \\
&= x \sum_{G \in \overline{\Se}(\M)}
  \bigl((x-1)\, \H_{\M|G} - \chi_{\M|G}\bigr)\,
  \H_{\M/(G \cup \{e\})}.
\end{align*}
The last equality is the defining convolution identity
$(x - 1)\, \H_L = \sum_{\emptyset \neq A \in \mathcal{L}(L)}
\chi_{L|A}\, \H_{L/A}$ for $L \neq \emptyset$.
Substitution in~\eqref{eq:chow-induction} yields
\[
(x-1)\, D_e(\M) = x(x-1) \sum_{G \in \overline{\Se}(\M)}
  \H_{\M|G}\, \H_{\M/(G \cup \{e\})},
\]
and cancellation of $x - 1$ proves~\eqref{eq:chow-del}.

It remains to treat a coloop~$e$.  Put $\N = \M/e = \M \setminus e$
and $J_L = \H_{L \oplus \U_{1,1}}$.
Here $\overline{\Se}(\M) = \mathcal{L}(\N) \setminus \{\emptyset\}$.
The case $\N = \emptyset$ is immediate.
For $\N \neq \emptyset$, pairing the flats $A$ and $A \cup \{e\}$ in
the defining convolution for~$\H$ gives
\[
(x-1)\, J_{\N}
= \sum_{\emptyset \neq A \in \mathcal{L}(\N)}
  \chi_{\N|A}\, J_{\N/A}
  + (x-1) \sum_{A \in \mathcal{L}(\N)}
  \chi_{\N|A}\, \H_{\N/A}.
\]
The second sum is $(x-1)\, x\, \H_{\N}$.
Applying the induction hypothesis to each $J_{\N/A}$ in the first sum
and reindexing once more gives, with
\[
C_{\N} = \sum_{\emptyset \neq B \in \mathcal{L}(\N)}
  \H_{\N|B}\, \H_{\N/B},
\]
\[
\sum_{\emptyset \neq A \in \mathcal{L}(\N)}
  \chi_{\N|A}\, J_{\N/A}
= (x-1)\, \H_{\N} + x\bigl((x-1)\, C_{\N} - (x-1)\, \H_{\N}\bigr).
\]
Indeed, the double sum arising from the induction hypothesis is reindexed
by $B = A \cup C$, and its inner sum is
\[
\sum_{\substack{A \in \mathcal{L}(\N|B) \\
    \emptyset \neq A \subsetneq B}}
  \chi_{\N|A}\, \H_{(\N|B)/A}
= (x-1)\, \H_{\N|B} - \chi_{\N|B}.
\]
Consequently,
\[
(x-1)\, J_{\N} = (x-1)(\H_{\N} + x\, C_{\N}),
\]
which is precisely~\eqref{eq:chow-del} for the coloop~$e$.
\end{proof}


\subsection{The Dual Chow Polynomial}
Let $\Hstar_{\N}(x)$ denote the dual Chow polynomial of a loopless
matroid~$\N$, with the convention $\Hstar_{\emptyset}(x) = 1$.
For the next lemma, we make the auxiliary convention that
$\Hstar_{\N}(x) = 0$ when $\N$ has a loop.
This does not change the convolution defining~$\Gstar_{\M}$ for
loopless~$\M$, because contractions by flats of a loopless matroid
are loopless.

\begin{lemma}[Loop-extended dual-Chow deletion]\label{lem:dual-chow-del}
Let $\N$ be a loopless matroid and let $p \in E(\N)$ be a non-coloop.
Then
\[
\Hstar_{\N}(x) = \Hstar_{\N \setminus p}(x) + \Hstar_{\N/p}(x)
  + x \sum_{K \in \overline{\Se[p]}(\N)}
    \Hstar_{\N|K}(x)\, \Hstar_{\N/(K \cup \{p\})}(x).
\]
\end{lemma}

\begin{proof}
If $p$ has no parallel mate, the deletion formula of
Caiolo--Ferroni--Hoster has the form
\[
\Hstar_{\N} = \Hstar_{\N \setminus p} + (1+x)\, \Hstar_{\N/p}
  + x \sum_{K \in \overline{\Se[p]}(\N)}
    \Hstar_{\N|K}\, \Hstar_{\N/(K \cup \{p\})}.
\]
In this case $\emptyset \in \Se[p](\N)$, and the $K = \emptyset$ summand
in the displayed formula of the lemma is $x\, \Hstar_{\N/p}$, because
$\Hstar_{\N|\emptyset} = \Hstar_{\emptyset} = 1$.
Thus the two forms agree.

Suppose instead that $p$ has a parallel mate~$q$.
Then $\overline{\Se[p]}(\N) = \emptyset$: a flat avoiding~$p$ also
avoids~$q$, whereas adjoining~$p$ without~$q$ cannot produce a flat.
Deleting~$p$ preserves the ranked lattice of flats, so
$\Hstar_{\N} = \Hstar_{\N \setminus p}$.
On the other hand, $q$~is a loop of~$\N/p$, and the auxiliary
zero-on-loops convention gives $\Hstar_{\N/p} = 0$.
This proves the stated recurrence in every case.
\end{proof}


\subsection{The $Z$-Polynomial}
The $Z$-polynomial of~$\M$ is defined by the convolution
\begin{equation}\label{eq:Z-def}
\ZP_{\M}(x) = (\zeta^{\rev} * \PP)_{\M}
= \sum_{A \in \mathcal{L}(\M)} x^{\rk(A)}\, \PP_{\M/A}(x),
\end{equation}
where $\zeta^{\rev}_{\N}(x) = x^{\rk(\N)}$.
The following is a direct specialization of \cref{cor:loop-vanish}.

\begin{corollary}[$Z$-polynomial deletion formula]\label{cor:Z-del}
Let $\M$ be loopless and let $e \in E(\M)$ be a non-coloop.  Then
\[
\ZP_{\M}(x) = \ZP_{\M \setminus e}(x)
  + \sum_{F \in \Se(\M)}
    \tau\bigl(\M/(F \cup \{e\})\bigr)\, x^{\crk_{\M}(F)/2}\,
    \ZP_{\M|F}(x).
\]
\end{corollary}

\begin{proof}
Extend $\PP$ by~$0$ to matroids with loops, and put
\[
\phi(Q) =
\begin{cases}
\tau(Q)\, x^{(\rk(Q)+1)/2}, & \text{if $Q$ is loopless and $\rk(Q)$ is odd},\\
0, & \text{otherwise}.
\end{cases}
\]
Apply \cref{cor:loop-vanish} with
$g_{\N} = x^{\rk(\N)}$ and $h_{\N} = \PP_{\N}(x)$.
The Kazhdan--Lusztig deletion formula gives
\[
\alpha_h = -x,\quad \Phi^h = \phi,\quad \Psi^h = 0,
\]
while $g_{\N} = g_{\N \setminus p}$ for a non-coloop~$p$ and
$g_{\N} = x\, g_{\N/p}$ for a coloop~$p$, so that
\[
\alpha_g = 0,\quad \gamma_g = x,\quad \Phi^g = \Psi^g = 0.
\]
The convolution in the Convolution Deletion Formula is
$f_{\N} = \ZP_{\N}(x)$, and its additional loop-vanishing hypotheses hold.
Since $\delta = -x + x = 0$, only the $\Phi^h$-term remains.
Finally, for $F \in \Se(\M)$,
\[
\phi\bigl(\M/(F \cup \{e\})\bigr)
= \tau\bigl(\M/(F \cup \{e\})\bigr)\, x^{\crk_{\M}(F)/2},
\]
because $\rk(\M/(F \cup \{e\})) + 1 = \crk_{\M}(F)$.
This proves the formula.
\end{proof}


\begin{lemma}\label{lem:Prev-del}
For any loopless matroid~$\M$ and $e \in E(\M)$ that is not a coloop,
the reversed Kazhdan--Lusztig polynomial satisfies:
\[
\PP^{\rev}_{\M}(x) = \PP^{\rev}_{\M \setminus e}(x)
  - \PP^{\rev}_{\M/e}(x)
  + \sum_{F \in \Se(\M)}
    \tau(\M/(F \cup \{e\}))\, x^{\crk_{\M}(F)/2}\,
    \PP^{\rev}_{\M|F}(x).
\]
\end{lemma}

\begin{proof}
The standard deletion formula for the Kazhdan--Lusztig polynomial
$\PP_{\M}(x)$ is:
\[
\PP_{\M}(x) = \PP_{\M \setminus e}(x) - x\, \PP_{\M/e}(x)
  + \sum_{F \in \Se(\M)}
    \tau(\M/(F \cup \{e\}))\, x^{\crk_{\M}(F)/2}\,
    \PP_{\M|F}(x).
\]
Substituting $x$ with $x^{-1}$ and multiplying by $x^{\rk(\M)}$
yields the stated formula for~$\PP^{\rev}_{\M}(x)$.
\end{proof}


\subsection{The Right Augmented Chow Polynomial}

\[
\FF_{\N}(x) = (\H * \PP^{\rev})_{\N}
= \sum_{A \in \mathcal{L}(\N)}
  \H_{\N|A}(x)\, \PP^{\rev}_{\N/A}(x)
\]
for every loopless matroid~$\N$.
For the bookkeeping required by the Convolution Deletion Formula,
extend~$\H$ by~$0$ to matroids with loops.

\begin{corollary}[Right augmented Chow deletion formula]\label{cor:aug-chow-del}
Let $\M$ be a simple matroid and let $e \in E(\M)$ be a non-coloop.
The right augmented Chow polynomial satisfies:
\begin{equation}\label{eq:aug-chow-del}
\FF_{\M}(x) = \FF_{\M \setminus e}(x)
  + x \sum_{G \in \overline{\Se}(\M)}
    \H_{\M|G}(x)\, \FF_{\M/(G \cup \{e\})}(x)
  + \sum_{F \in \Se(\M)}
    \tau(\M/(F \cup \{e\}))\, x^{\crk_{\M}(F)/2}\,
    \FF_{\M|F}(x).
\end{equation}
\end{corollary}

\begin{proof}
Extend $\PP^{\rev}$ by~$0$ to matroids with loops, and set
\[
\phi(Q) =
\begin{cases}
\tau(Q)\, x^{(\rk(Q)+1)/2}, & \text{if $Q$ is loopless and $\rk(Q)$ is odd},\\
0, & \text{otherwise}.
\end{cases}
\]
We apply \cref{cor:loop-vanish} with
\[
g_{\N} = \H_{\N}(x),\quad h_{\N} = \PP^{\rev}_{\N}(x).
\]
\Cref{cor:chow-del} gives the required $g$-recurrences with
\[
\alpha_g = 0,\quad \gamma_g = 1,\quad \Psi^g = 0,
\]
and
\[
\Phi^g(Q) =
\begin{cases}
x\, \H_Q(x), & \text{if } E(Q) \neq \emptyset, \\
0, & \text{if } Q = \emptyset.
\end{cases}
\]
Indeed, the non-coloop formula has precisely the $\overline{\Se[p]}(\N)$
sum, while the coloop formula has the extra term $\H_{\N/p}$ required
by the Convolution Deletion Formula.
\Cref{lem:Prev-del} supplies the $h$-recurrence with
\[
\alpha_h = -1,\quad \Phi^h = \phi,\quad \Psi^h = 0.
\]
The equality
\[
\rk\bigl(\N/(K \cup \{p\})\bigr) + 1 = \crk_{\N}(K)
\]
for $K \in \Se[p](\N)$ identifies its $\Phi^h$-weight with the boundary
coefficient in that lemma.

The convolution in the Convolution Deletion Formula is $f_{\N} = \FF_{\N}$,
and $h$ vanishes on looped matroids.  The remaining hypotheses of
\cref{cor:loop-vanish} hold because $\Psi^h = 0$ and $\Phi^g(\emptyset) = 0$.
Moreover,
\[
\delta = \alpha_h + \gamma_g - \alpha_g = -1 + 1 = 0.
\]
Consequently, the $\Phi^h$-sum in the loop-vanishing corollary is
\[
\sum_{F \in \Se(\M)}
  \tau\bigl(\M/(F \cup \{e\})\bigr)\, x^{\crk_{\M}(F)/2}\,
  \FF_{\M|F}(x),
\]
whereas its $\Phi^g$-sum is
\[
x \sum_{G \in \overline{\Se}(\M)}
  \H_{\M|G}(x)\, \FF_{\M/(G \cup \{e\})}(x).
\]
These are exactly the two terms in the stated deletion formula.
\end{proof}


\subsection{The Left Augmented Dual Chow Polynomial}
The left augmented dual Chow polynomial is defined by
\[
\Gstar_{\M}(x) = \sum_{F \in \mathcal{L}(\M)}
  \QQ^{\rev}_{\M|F}(x)\, \Hstar_{\M/F}(x),
\]
where $\Hstar_{\M}(x)$ is the dual Chow polynomial and
$\QQ^{\rev}_{\M}(x) = x^{\rk(\M)}\, \QQ_{\M}(x^{-1})$ is the
reversed inverse Kazhdan--Lusztig polynomial.

\begin{corollary}[Left augmented dual Chow deletion formula]
\label{cor:Gstar-del}
Let $\M$ be a simple matroid, and let $e \in E(\M)$ be a non-coloop.
Then the left augmented dual Chow polynomial $\Gstar_{\M}(x)$ satisfies
the deletion formula:
\begin{align*}
\Gstar_{\M}(x) &= \Gstar_{\M \setminus e}(x)
  + (1+x)\, \Gstar_{\M/e}(x)
  + x \sum_{F \in \overline{\Se}(\M)}
    \Gstar_{\M|F}(x)\, \Hstar_{\M/(F \cup \{e\})}(x) \\
&\quad - \sum_{\substack{K \in \Te(\M) \\ \rk_{\M}(K) \text{ even}}}
    \tau\bigl((\M|K)/e\bigr)\, x^{\rk_{\M}(K)/2}\,
    \Gstar_{\M/K}(x).
\end{align*}
\end{corollary}

\begin{proof}
Put $q_{\N} = \QQ^{\rev}_{\N}(x)$, using the zero-on-loops convention
already fixed for~$\QQ$, and define
\[
\psi(Q) =
\begin{cases}
-\tau(Q)\, x^{(\rk(Q)+1)/2}, & \text{if $Q$ is loopless and $\rk(Q)$ is odd},\\
0, & \text{otherwise}.
\end{cases}
\]
The inverse Kazhdan--Lusztig deletion formula, proved below as
\cref{cor:invKL-del} gives, after reversal, for a loopless pointed
matroid $(\N, p)$ with $p$ a non-coloop,
\[
q_{\N} = q_{\N \setminus p} + (1+x)\, q_{\N/p}
  + \sum_{K \in \Te[p](\N)}
    \psi\bigl((\N|K)/p\bigr)\, q_{\N/K}.
\]
If $p$ is a coloop, then $\QQ_{\N} = \QQ_{\N/p}$, and reversal gives
\[
q_{\N} = x\, q_{\N/p}.
\]
We may therefore apply \cref{cor:loop-vanish} with
\[
h_{\N} = \Hstar_{\N}(x),\quad g_{\N} = q_{\N},
\]
and parameters
\[
\alpha_h = 1,\quad
\Phi^h(Q) = x\, \Hstar_Q(x),\quad
\Psi^h = 0,\quad
\alpha_g = 1+x,\quad
\gamma_g = x,\quad
\Phi^g = 0,\quad
\Psi^g(Q) = \psi(Q).
\]
\Cref{lem:dual-chow-del} supplies the $h$-recurrence, including the case
in which $p$ has a parallel mate.  Its auxiliary extension also gives the
required loop-vanishing condition.
The convolution in the Convolution Deletion Formula is
\[
f_{\N} = \sum_{A \in \mathcal{L}(\N)}
  \QQ^{\rev}_{\N|A}(x)\, \Hstar_{\N/A}(x) = \Gstar_{\N}(x),
\]
and
\[
\delta = \alpha_h + \gamma_g - \alpha_g = 1 + x - (1+x) = 0.
\]
Thus the $\Phi^h$-term in the loop-vanishing corollary is
\[
x \sum_{F \in \overline{\Se}(\M)}
  \Gstar_{\M|F}(x)\, \Hstar_{\M/(F \cup \{e\})}(x),
\]
and the $\Psi^g$-term is
\[
\sum_{K \in \Te(\M)}
  \psi\bigl((\M|K)/e\bigr)\, \Gstar_{\M/K}(x).
\]
Since $\rk\bigl((\M|K)/e\bigr) = \rk_{\M}(K) - 1$, the latter sum is
exactly
\[
-\sum_{\substack{K \in \Te(\M) \\ \rk_{\M}(K) \text{ even}}}
  \tau\bigl((\M|K)/e\bigr)\, x^{\rk_{\M}(K)/2}\,
  \Gstar_{\M/K}(x).
\]
Together with the deletion and contraction terms of the master corollary,
this is the displayed formula.
\end{proof}

\begin{remark}[The $\Gstar$ contraction coefficient]\label{rem:Gstar-coeff}
For the right augmented Chow polynomial, the nonzero master boundary
weight is~$\Phi^g$, so its extra contribution is indexed by~$\Se(\M)$.
Here $\Phi^g = 0$; instead, the $\Gstar$ specialization has the
contraction coefficient
$\Psi^g(Q) = \psi(Q)$.
Consequently the Convolution Deletion Formula evaluates it
on the contraction $(\M|K)/e$ and produces
$\Psi^g\bigl((\M|K)/e\bigr) = \psi\bigl((\M|K)/e\bigr)$.
Thus the contraction by~$e$ is taken after restricting to~$K$,
which is the source of the $\Te(\M)$-sum in the $\Gstar$ formula.
\end{remark}


\subsection{The Inverse Kazhdan--Lusztig Polynomial}
For a loopless matroid~$\M$, put
\[
\hat{\QQ}_{\M}(x) = (-1)^{\rk(\M)}\, \QQ_{\M}(x).
\]
Extend $\PP$, $\QQ$, and~$\tau$ by zero to matroids with loops.
With this convention, $\hat{\QQ}$ is the right convolution inverse of~$\PP$:
\begin{equation}\label{eq:PQ-inv}
\sum_{A \in \mathcal{L}(\M)}
  \PP_{\M|A}(x)\, \hat{\QQ}_{\M/A}(x) = \delta_{\emptyset, \M}.
\end{equation}
For looped~$\M$, this identity is automatic, because every flat
contains a loop and every restricted $\PP$-factor therefore vanishes.
\Cref{cor:inv-loop-vanish} therefore gives the deletion formula
formally from the deletion formula for~$\PP$.

\begin{corollary}[Inverse Kazhdan--Lusztig deletion formula]
\label{cor:invKL-del}
Let $\M$ be a loopless matroid and let $e \in E(\M)$ be a non-coloop.
Then
\begin{equation}\label{eq:invKL-del}
\QQ_{\M}(x) = \QQ_{\M \setminus e}(x)
  + (1+x)\, \QQ_{\M/e}(x)
  - \sum_{\substack{F \in \Te(\M) \\ \rk(F) \text{ even}}}
    \tau\bigl((\M|F)/e\bigr)\, x^{\rk(F)/2}\,
    \QQ_{\M/F}(x).
\end{equation}
\end{corollary}

\begin{proof}
Since $\PP_{\emptyset} = 1$, \cref{cor:inv-loop-vanish} applies with
\[
\kappa_{\N} = \PP_{\N}(x),\quad
\lambda_{\N} = \hat{\QQ}_{\N}(x),\quad
\alpha = -x,\quad
\gamma = 1,
\]
and define
\[
\phi(\N) =
\begin{cases}
\tau(\N)\, x^{(\rk(\N)+1)/2},
  & \text{if $\N$ is loopless and $\rk(\N)$ is odd},\\
0, & \text{otherwise}.
\end{cases}
\]
Indeed, the Kazhdan--Lusztig deletion recurrence has the form
\[
\PP_{\N} = \PP_{\N \setminus p} - x\, \PP_{\N/p}
  + \sum_{G \in \Se[p](\N)}
    \phi\bigl(\N/(G \cup \{p\})\bigr)\, \PP_{\N|G},
\]
because $\rk\bigl(\N/(G \cup \{p\})\bigr) + 1 = \crk_{\N}(G)$.
Its coloop relation is $\PP_{\N} = \PP_{\N/p}$.
Thus the corollary gives
\[
\hat{\QQ}_{\M} = \hat{\QQ}_{\M \setminus e}
  - (1+x)\, \hat{\QQ}_{\M/e}
  - \sum_{\substack{F \in \Te(\M) \\ \rk(F) \text{ even}}}
    \tau\bigl((\M|F)/e\bigr)\, x^{\rk(F)/2}\,
    \hat{\QQ}_{\M/F}.
\]
The rank of $(\M|F)/e$ is $\rk(F) - 1$, so only even-rank~$F$
can contribute.  If $r = \rk(\M)$, then
\[
\rk(\M \setminus e) = r,\quad
\rk(\M/e) = r-1,\quad
\rk(\M/F) = r - \rk(F).
\]
Replacing $\hat{\QQ}_{\N}$ by $(-1)^{\rk(\N)}\, \QQ_{\N}$ therefore
reverses the contraction sign and preserves every boundary sign,
proving~\eqref{eq:invKL-del}.
\end{proof}


\subsection{The Inverse $Z$-Polynomial}
For a loopless matroid~$\M$, put
\[
\hat{\YY}_{\M}(x) = (-1)^{\rk(\M)}\, \YY_{\M}(x).
\]
For the following convolutional argument, extend $\ZP$, $\hat{\YY}$
(and hence~$\YY$), and~$\tau$ by zero to matroids with loops.
This is an auxiliary bookkeeping convention; it does not
evaluate~\eqref{eq:Z-def} on a looped matroid.
With this convention, $\hat{\YY}$ is the right convolution inverse
of~$\ZP$:
\begin{equation}\label{eq:ZY-inv}
\sum_{A \in \mathcal{L}(\M)}
  \ZP_{\M|A}(x)\, \hat{\YY}_{\M/A}(x) = \delta_{\emptyset, \M}.
\end{equation}
For looped~$\M$, this identity is automatic, because every flat
contains a loop and every restricted $\ZP$-factor therefore vanishes.

\begin{corollary}[Inverse $Z$-polynomial deletion formula]
\label{cor:invZ-del}
Let $\M$ be a loopless matroid and let $e \in E(\M)$ be a non-coloop.
Then
\begin{equation}\label{eq:invZ-del}
\YY_{\M}(x) = \YY_{\M \setminus e}(x)
  + (1+x)\, \YY_{\M/e}(x)
  - \sum_{\substack{F \in \Te(\M) \\ \rk(F) \text{ even}}}
    \tau\bigl((\M|F)/e\bigr)\, x^{\rk(F)/2}\,
    \YY_{\M/F}(x).
\end{equation}
\end{corollary}

\begin{proof}
Since $\ZP_{\emptyset} = 1$, \cref{cor:inv-loop-vanish} applies with
\[
\kappa_{\N} = \ZP_{\N}(x),\quad
\lambda_{\N} = \hat{\YY}_{\N}(x),\quad
\alpha = 0,\quad
\gamma = 1+x,
\]
and
\[
\phi(\N) =
\begin{cases}
\tau(\N)\, x^{(\rk(\N)+1)/2},
  & \text{if $\N$ is loopless and $\rk(\N)$ is odd},\\
0, & \text{otherwise}.
\end{cases}
\]
The loopless deletion formula of \cref{cor:Z-del} has the required form
with $\alpha = 0$, since
$\rk\bigl(\N/(G \cup \{p\})\bigr) + 1 = \crk_{\N}(G)$.
If $p$ is a coloop of~$\N$, the flats of~$\N$ occur in pairs $A$
and $A \cup \{p\}$ with $A \in \mathcal{L}(\N/p)$.
Adjoining a coloop does not change the Kazhdan--Lusztig polynomial,
so the corresponding summands in~\eqref{eq:Z-def} are $1$ and~$x$
times the same summand of~$\ZP_{\N/p}$.  Hence
\[
\ZP_{\N} = (1+x)\, \ZP_{\N/p}.
\]
Thus $\alpha - \gamma = -(1+x)$, and the corollary gives
\[
\hat{\YY}_{\M} = \hat{\YY}_{\M \setminus e}
  - (1+x)\, \hat{\YY}_{\M/e}
  - \sum_{\substack{F \in \Te(\M) \\ \rk(F) \text{ even}}}
    \tau\bigl((\M|F)/e\bigr)\, x^{\rk(F)/2}\,
    \hat{\YY}_{\M/F}.
\]
Again, a nonzero boundary term has $\rk(F)$ even.
If $r = \rk(\M)$, then
\[
\rk(\M \setminus e) = r,\quad
\rk(\M/e) = r-1,\quad
\rk(\M/F) = r - \rk(F).
\]
Replacing $\hat{\YY}_{\N}$ by $(-1)^{\rk(\N)}\, \YY_{\N}$ therefore
reverses the contraction sign and preserves every boundary sign,
proving~\eqref{eq:invZ-del}.
\end{proof}

\begin{remark}[Why the inverse formulas coincide]\label{rem:inverse-coincide}
The deletion formulas in \cref{cor:invKL-del,cor:invZ-del} coincide
because the forward kernels have the same boundary weight~$\phi$, and
\[
(-x) - 1 = -(1+x) = 0 - (1+x).
\]
Equivalently, the pairs $(\alpha, \gamma) = (-x, 1)$ for~$\PP$
and $(0, 1+x)$ for~$\ZP$ have the same difference $\alpha - \gamma$.
The loop-vanishing inverse formula depends on the forward recurrence
precisely through this difference and~$\phi$; after the same rank-sign
normalization, it therefore yields identical recurrences for~$\QQ$
and~$\YY$.
\end{remark}


%%% ============================================================
%%% SECTION 6 -- Remark on Gamma-Positivity
%%% ============================================================
\section{Remark on Gamma-Positivity}

\begin{remark}\label{rem:gamma-pos}
Recall that a symmetric polynomial $p(x)$ of degree~$d$ is
\emph{$\gamma$-positive} if it is a nonnegative linear combination of
the polynomials $x^i(1+x)^{d-2i}$.
The deletion formula for the right augmented Chow polynomial is
compatible with its symmetry, but it does not by itself imply
$\gamma$-positivity: its summands have different symmetry degrees,
and the formula includes the coefficients~$\tau$ of Kazhdan--Lusztig
polynomials.  Establishing $\gamma$-positivity therefore requires an
additional structural argument, which is not supplied here.
\end{remark}


%%% ============================================================
%%% SECTION 7 -- Examples and Specializations
%%% ============================================================
\section{Examples and Specializations}

We illustrate the deletion formula for the right augmented Chow
polynomial in several matroid classes.  The following calculations
specialize the formula proved above.  We write
\[
\B_j := \U_{j,j},\qquad \C_j := \U_{j-1,j},
\]
for the Boolean matroid of rank~$j$ and the circuit of size~$j$,
respectively.  We set
\[
A_j(x) := \H_{\B_j}(x),\qquad
\widetilde{A}_j(x) := \FF_{\B_j}(x),
\]
which denote the Chow and right augmented Chow polynomials of the
Boolean matroid of rank~$j$.

The class-specific calculations in the subsequent sections specialize
\cref{cor:aug-chow-del} for a simple matroid~$\M$:
\begin{equation}\tag{D$_{\FF}$}\label{eq:DF}
\FF_{\M}(x) = \FF_{\M \setminus e}(x)
  + x \sum_{G \in \overline{\Se}(\M)}
    \H_{\M|G}(x)\, \FF_{\M/(G \cup \{e\})}(x)
  + \sum_{G \in \Se(\M)}
    \tau(\M/(G \cup \{e\}))\, x^{\crk_{\M}(G)/2}\,
    \FF_{\M|G}(x).
\end{equation}
Whenever a Kazhdan--Lusztig $\tau$-coefficient $\tau(\N)$ is nonzero,
$\rk(\N)$ is odd; hence the corresponding half-integral exponent
in~\eqref{eq:DF} is an integer.


\subsection{Uniform Matroids}

\begin{theorem}\label{thm:uniform-del}
Let $2 \le r < n$, let $\M = \U_{r,n}$, and let $e \in E(\M)$.  Then
\[
\Se(\M) = \{G \subseteq E(\M) \setminus \{e\} : |G| \le r-2\}.
\]
Consequently, the augmented Chow polynomial satisfies
\begin{equation}\label{eq:uniform-del}
\FF_{\U_{r,n}}(x) = \FF_{\U_{r,n-1}}(x)
  + x \sum_{j=1}^{r-2} \binom{n-1}{j}\,
    A_j(x)\, \FF_{\U_{r-j-1,n-j-1}}(x)
  + \sum_{j=0}^{r-2} \binom{n-1}{j}\,
    \tau(\U_{r-j-1,n-j-1})\, x^{(r-j)/2}\,
    \widetilde{A}_j(x).
\end{equation}
The first sum is empty when $r = 2$.
\end{theorem}

\begin{proof}
The proper flats of~$\U_{r,n}$ are precisely the subsets of cardinality
strictly less than~$r$.  Hence, for a flat~$G$ not containing~$e$,
$G \cup \{e\} \in \mathcal{L}(\U_{r,n})$ if and only if $|G| + 1 < r$.
Therefore,
\[
G \in \Se(\M) \iff |G| \le r-2.
\]
If $|G| = j \le r-2$, then $G$ is independent, so
$\M|G \cong \B_j$ and $\M/(G \cup \{e\}) \cong \U_{r-j-1,n-j-1}$.
Moreover, $\crk_{\M}(G) = r - j$.
There are exactly $\binom{n-1}{j}$ such flats.
Thus, the quadratic term in~\eqref{eq:DF} becomes
\[
x \sum_{j=1}^{r-2} \binom{n-1}{j}\,
  A_j(x)\, \FF_{\U_{r-j-1,n-j-1}}(x),
\]
while the $\tau$-term becomes
\[
\sum_{j=0}^{r-2} \binom{n-1}{j}\,
  \tau(\U_{r-j-1,n-j-1})\, x^{(r-j)/2}\,
  \widetilde{A}_j(x).
\]
Finally, $\M \setminus e \cong \U_{r,n-1}$, which
establishes~\eqref{eq:uniform-del}.
\end{proof}

\begin{remark}[Boundary Cases]\label{rem:uniform-boundary}
For $r = 2$, we have
$\FF_{\U_{2,n}}(x) = \FF_{\U_{2,n-1}}(x) + x$,
because $\tau(\U_{1,n-1}) = 1$.
The case $r = n$ is excluded because every element is a coloop.
The case $r = 1 < n$ is excluded from the stated simple-matroid
theorem because $\U_{1,n}$ is not simple.
\end{remark}


\subsection{Glued Cycles}
Let $\mathcal{C}_{a,b}$ be the graphic matroid of two cycles of
lengths~$a$ and~$b$, with $a,b \ge 3$, glued along a common edge~$e$.
Put $p := a-1$ and $q := b-1$.
The graph consists of~$e$ and two internally disjoint paths of
lengths~$p$ and~$q$ between the endpoints of~$e$.

\begin{lemma}\label{lem:glued-Se}
A subset $G \subseteq E(\mathcal{C}_{a,b}) \setminus \{e\}$ lies in
$\Se(\mathcal{C}_{a,b})$ if and only if $G$ omits at least two edges
from each of the two paths.
\end{lemma}

\begin{proof}
Let $P_a$ and~$P_b$ denote the two paths.
If $G$ omits only one edge~$f$ from~$P_a$, then in $G \cup \{e\}$
the endpoints of~$f$ are connected by the remaining edges of~$P_a$
together with~$e$.
Thus $f \in \cl(G \cup \{e\})$, so $G \cup \{e\}$ is not a flat.
The same applies to~$P_b$.

Conversely, suppose at least two edges are omitted from each path.
Then neither~$G$ nor $G \cup \{e\}$ contains all but one edge of any
cycle of the theta graph.
Hence, no omitted edge has both endpoints in the same connected
component of the corresponding spanning subgraph.
Therefore, both~$G$ and $G \cup \{e\}$ are flats.
\end{proof}

\begin{lemma}\label{lem:glued-minors}
Suppose $G \in \Se(\mathcal{C}_{a,b})$ omits $s$~edges from~$P_a$
and $t$~edges from~$P_b$.  Then $s \ge 2$, $t \ge 2$, and
\[
\mathcal{C}_{a,b}|G \cong \B_{p+q-s-t},\qquad
\mathcal{C}_{a,b}/(G \cup \{e\}) \cong \C_s \oplus \C_t.
\]
\end{lemma}

\begin{proof}
The restriction to~$G$ is a forest, so every element of~$G$ is a
coloop in the restriction.  Hence it is Boolean.
After contracting~$e$ and all selected edges of each path,
the $s$~uncontracted edges of the first path form a cycle, as do
the $t$~uncontracted edges of the second path.
These two cycles meet only at a vertex, so their cycle matroids
form a direct sum.
\end{proof}

\begin{lemma}\label{lem:glued-tau}
For $s, t \ge 2$, $\tau(\C_s \oplus \C_t) = 0$.
\end{lemma}

\begin{proof}
The Kazhdan--Lusztig polynomial is multiplicative under direct sums:
$\PP_{\C_s \oplus \C_t}(x) = \PP_{\C_s}(x)\, \PP_{\C_t}(x)$.
If $(s-1) + (t-1)$ is even, the assertion follows from the definition
of~$\tau$.  If it is odd, precisely one of $s-1$ and $t-1$ is even.
The degree bounds $\deg \PP_{\C_s}(x) < \frac{s-1}{2}$ and
$\deg \PP_{\C_t}(x) < \frac{t-1}{2}$ imply
$\deg \PP_{\C_s \oplus \C_t}(x) < \frac{s+t-3}{2}$, so the
coefficient defining~$\tau$ vanishes.
\end{proof}

\begin{theorem}\label{thm:glued-del}
The deletion formula for the common edge is purely quadratic:
\begin{equation}\label{eq:glued-del}
\FF_{\mathcal{C}_{a,b}}(x) = \FF_{\C_{a+b-2}}(x)
  + x \sum_{\substack{s=2 \\ t=2 \\
      (s,t) \neq (a-1,b-1)}}^{\substack{a-1 \\ b-1}}
    \binom{a-1}{s} \binom{b-1}{t}\,
    A_{a+b-2-s-t}(x)\, \FF_{\C_s \oplus \C_t}(x).
\end{equation}
\end{theorem}

\begin{proof}
Deleting the common edge joins the two paths into a single cycle of
length $(a-1) + (b-1) = a+b-2$.  Thus
$\mathcal{C}_{a,b} \setminus e \cong \C_{a+b-2}$.
For fixed~$s$ and~$t$, the omitted edges can be chosen in
$\binom{a-1}{s} \binom{b-1}{t}$ ways.
\Cref{lem:glued-minors} identifies the restriction and contraction
appearing in~\eqref{eq:DF}.
The pair $(s,t) = (a-1,b-1)$ corresponds to $G = \emptyset$, which
must be omitted from the quadratic sum.
\Cref{lem:glued-tau} dictates that the entire $\tau$-sum identically
vanishes.
\end{proof}


\subsection{Projective Geometry Deletions}
Let $\M = \PG(r-1, q)$ with $r \ge 2$, and let $e$ be a point.

\begin{lemma}\label{lem:PG-Se}
For a projective geometry, $\overline{\Se}(\M) = \{\emptyset\}$.
\end{lemma}

\begin{proof}
Both $\emptyset$ and $\{e\}$ are flats, so $\emptyset \in \Se(\M)$.
Suppose $G \neq \emptyset$ is a flat not containing~$e$.
Choose a point $f \in G$.
Since $G$ is a projective subspace not containing~$e$, its intersection
with the projective line $\langle e, f \rangle$ is exactly~$\{f\}$.
The line has at least three points, so
$\langle e, f \rangle \setminus \{e, f\} \neq \emptyset$.
All these additional points lie in the closure of $G \cup \{e\}$, but
not in $G \cup \{e\}$.  Therefore, $G \cup \{e\}$ is not a flat.
\end{proof}

\begin{lemma}\label{lem:PG-tau}
For $\M = \PG(r-1, q)$,
\[
\tau(\M/e) =
\begin{cases}
1, & r = 2, \\
0, & r \ge 3.
\end{cases}
\]
\end{lemma}

\begin{proof}
Projective geometries are modular, meaning $\PP_{\M/e}(x) = 1$.
If $r = 2$, $\M/e$ has rank~$1$, and consequently
$\tau(\M/e) = [x^0]\, 1 = 1$.
If $r \ge 3$, either $r-1$ is even (so $\tau$ vanishes by definition),
or $r-1$ is odd and the coefficient defining~$\tau$ has positive degree
in the constant polynomial~$1$, yielding zero.
\end{proof}

\begin{theorem}\label{thm:PG-del}
A single-point deletion satisfies
\begin{equation}\label{eq:PG-del}
\FF_{\PG(r-1,q) \setminus e}(x)
= \FF_{\PG(r-1,q)}(x) - \mathbf{1}_{\{r=2\}}\, x.
\end{equation}
In particular,
$\FF_{\PG(r-1,q) \setminus e}(x) = \FF_{\PG(r-1,q)}(x)$
for every $r \ge 3$.
\end{theorem}

\begin{proof}
Since $\overline{\Se}(\M) = \emptyset$, the quadratic sum
in~\eqref{eq:DF} is empty.  The only term in the $\tau$-sum is
$G = \emptyset$, for which $\crk_{\M}(\emptyset) = r$ and
$\FF_{\M|\emptyset}(x) = 1$.
Thus $\FF_{\M}(x) = \FF_{\M \setminus e}(x) + \tau(\M/e)\, x^{r/2}$,
and \cref{lem:PG-tau} yields the result.
\end{proof}


%%% ============================================================
%%% References
%%% ============================================================
%% Uncomment when ready:
%\printbibliography


\end{document}
